\hypertarget{appendix-46da295c6787}{}
\section*{Covectors linear functionals and normal transforms}
\section{Covectors linear functionals and normal transforms}\label{covectors-linear-functionals-and-normal-transforms}

\textbf{Curriculum depth:} Developed · \textbf{First seen:} Phase6 · \textbf{Latest development:} Phase7

\subsection{Definition and mechanism}\label{definition-and-mechanism}

Phase 7 derives \((RS)^{-T}=RS^{-1}\) for orthogonal \(R\) and nonzero diagonal scales \(S\), explains why zero scale invalidates inverse-based normal transport, and distinguishes shading-normal direction from determinant-signed cross-derived area. \href{Continuous\%20Notes/Computing\%20Foundations\%20-\%20From\%20Bits\%20to\%20C++\%20and\%20Systems\%20APIs\#Normals\%20undo\%20scale\%20before\%20following\%20rotation}{TRS normal derivation}.

A covector is a linear scalar-valued measurement rule. Surface normals represent rules annihilating tangent movement. Under an invertible active linear transform, preserved measurements require inverse-transpose normal transformation. A metric identifies vector and covector representations.

For tangent \((1,1,0)\) and normal \((1,-1,0)\), doubling horizontal scale requires normal \((1/2,-1,0)\) to preserve perpendicularity.

\subsection{Prerequisites}\label{prerequisites}

\begin{itemize}
\tightlist
\item
  \hyperref[appendix-ce0341f7ea0d]{Basis coordinates and dimension} --- An ordered basis fixes coordinate meaning and the space whose geometry is measured.
\item
  \href{}{Dot products and orthogonality (02-Foundations-II, appendix-8bc4e8fc547c)} --- Inner-product measurements supply the perpendicularity and alignment constraints.
\end{itemize}

\subsection{Used by}\label{used-by}

Later curriculum sources can deepen this same node; no future coverage is marked complete.

\subsection{Recall and next depth}\label{recall-and-next-depth}

Derive inverse transpose, distinguish active/passive changes, and handle singular transformations explicitly. Full production proofs, target-specific integration and learner mastery require further evidence.

\subsection{Source and navigation}\label{source-and-navigation}

\href{Continuous\%20Notes/Computing\%20Foundations\%20-\%20From\%20Bits\%20to\%20C++\%20and\%20Systems\%20APIs\#A\%20normal\%20is\%20a\%20measuring\%20rule\%20before\%20it\%20is\%20an\%20arrow}{Teaching narrative} · \hyperref[appendix-fd090d982a15]{Coverage and source gap} · \hyperref[appendix-a103e7846d21]{Evidence and limits} · \href{}{Dependency Map (02-Foundations-II, appendix-4ad47bdd269a)}
